\begin{helper}[Pushforward Agreement]
Fix $n\in\mathbb N$, $I\subseteq\mathrm{Fin}(n)$, parameter $\varepsilon\in\mathbb R$, parameter $M\in\mathbb N$, and $p_0\in\mathbb R$.
Let $\pi:\mathrm{Fin}(n)\hookrightarrow\mathrm{Fin}(M)\times\mathrm{Fin}(M)$ be an embedding.
Let $Z_{\mathrm{prod}}(v)=\noderef{FixedSetProductModelVar}(I,\varepsilon,p_0,\pi,v)$.
Let $Z_I(\omega)$ be the sample-based quantity defined identically as in \(\noderef{FixedSetProductModelTail}\).
Let $P_R=\{\pi(u)_1:u\in I\}$, $P_B=\{\pi(u)_2:u\in I\}$.
Let $D(\omega)$ denote the incident data extracted from sample $\omega$, and let $D(v)$ denote the incident data extracted from product assignment $v$, defined identically as in \(\noderef{FixedSetProductModelTailPushforwardGraphFiberSum}\) and \(\noderef{FixedSetProductModelTailPushforwardProductFiberSum}\).
For any predicate $H$, define the induced data predicate $H_D(D)$ as $\exists v: D(v)=D \land H(Z_{\mathrm{prod}}(v))$.
Then for any sample $\omega$ with \(\noderef{TwoBiteEmbedding}(\omega)=\pi\), we have the logical equivalence
\[
H(Z_I(\omega)) \iff H_D(D(\omega)),
\]
and as sets of product assignments, we have the equality
\[
\{v: H(Z_{\mathrm{prod}}(v))\} = \{v: H_D(D(v))\}.
\]
\end{helper}
\begin{proof}
Fix a sample \(\omega\) in the cell, so
\(\noderef{TwoBiteEmbedding}(\omega)=\pi\).  Define a product assignment
\(v_\omega\) from the incident data of \(\omega\) as follows.  For a red
coordinate \(x\), set the first component of \(v_{\omega,x}\) to
\[
D_R(\omega)(x)=
\begin{cases}
\emptyset, & x\in P_R,\\
\{r\in P_R:\text{\(x\) is adjacent to \(r\) in the red graph of \(\omega\)}\},
  & x\notin P_R,
\end{cases}
\]
and set the redundant second component arbitrarily, say to \(\emptyset\).
For a blue coordinate \(y+M\), set the second component to the analogous set
\(D_B(\omega)(y)\), and set the redundant first component to \(\emptyset\).
Then, by construction, \(D(v_\omega)=D(\omega)\).

We now check that this equality of incident data forces equality of the two
values.  In the sample \(\omega\), because the embedding is \(\pi\), every
vertex \(u\in I\) has red projection \(\pi(u)_1\in P_R\) and blue projection
\(\pi(u)_2\in P_B\).  In the product sample defining
\(\noderef{FixedSetProductModelVar}\), the red graph is
\[
\noderef{FixedSetFakeGraph}
  \bigl(P_R,\;x\mapsto (v_x)_1\bigr),
\]
and the blue graph is the corresponding fake graph on \(P_B\) using the
second components of the last \(M\) coordinates.

For a red base vertex \(x\notin P_R\) and \(u\in I\), the red lifted
neighborhood condition in \(\omega\) is exactly the assertion that
\(\pi(u)_1\in D_R(\omega)(x)\).  The red lifted neighborhood condition in the
product sample is, after unfolding \(\noderef{FixedSetFakeGraph}\), exactly
the assertion that \(\pi(u)_1\in (v_x)_1\).  If \(D(v)=D(\omega)\), these two
conditions are identical.  If \(x\in P_R\), then the base vertex
\(\operatorname{inl}x\) belongs to the embedded set removed before the
external-small-class stage, so its neighborhood value is irrelevant to the
final count.  The same argument applies to blue base vertices \(y\notin P_B\),
using \(D_B\) and the second components of the blue coordinates.

It follows that for every external red or blue base vertex \(x\), the set
\(\noderef{TwoBiteX}(\omega,I,x)\) equals the corresponding
\(\noderef{TwoBiteX}\) set in the product sample.  Therefore their
cardinalities are equal, so the \(\noderef{TwoBiteSmallClass}\) predicate
selects exactly the same external base vertices.  For each selected base
vertex, the input set to \(\noderef{TwoBiteFinalPairs}\) is the same, hence
the final-pair set is the same.  The red and blue projection maps used in the
projection-distinct filter are also the same, because both samples have
embedding \(\pi\).  Thus the union of final pairs and its
projection-distinct subfilter are equal, so their cardinalities are equal:
\[
Z_I(\omega)=Z_{\mathrm{prod}}(v)
\quad\text{whenever }D(v)=D(\omega).
\]
In particular \(Z_I(\omega)=Z_{\mathrm{prod}}(v_\omega)\).

This proves the first equivalence.  If \(H(Z_I(\omega))\) holds, the
assignment \(v_\omega\) satisfies \(D(v_\omega)=D(\omega)\) and
\[
H(Z_{\mathrm{prod}}(v_\omega))=H(Z_I(\omega)),
\]
so \(H_D(D(\omega))\) holds.  Conversely, if \(H_D(D(\omega))\) holds, choose
\(v\) with \(D(v)=D(\omega)\) and \(H(Z_{\mathrm{prod}}(v))\).  The equality
just proved gives \(Z_{\mathrm{prod}}(v)=Z_I(\omega)\), hence
\(H(Z_I(\omega))\).

For the product-assignment set equality, one inclusion is immediate: if
\(H(Z_{\mathrm{prod}}(v))\), take \(v'=v\) in the definition of
\(H_D(D(v))\).  Conversely, if \(H_D(D(v))\), choose \(v'\) with
\(D(v')=D(v)\) and \(H(Z_{\mathrm{prod}}(v'))\).  The same incident-data
argument applied to the two product samples gives
\[
Z_{\mathrm{prod}}(v')=Z_{\mathrm{prod}}(v),
\]
because their fake red and blue graphs have the same adjacencies relevant to
all \(\noderef{TwoBiteX}\), \(\noderef{TwoBiteSmallClass}\),
\(\noderef{TwoBiteFinalPairs}\), and projection-distinct filters.  Hence
\(H(Z_{\mathrm{prod}}(v))\).  The two filtered sets of product assignments are
therefore equal.
\end{proof}
